In this chapter, we present the shortest superstring approximation algorithms of
Breslauer, Jiang, and Jiang \cite{Breslauer1996}. We first describe the generic
approximation algorithm, then establish the overlap-rotation lemma, the main
theorem on which the approximation guarantees of both algorithms rely, and
finally present the two algorithms themselves, achieving approximation ratios of
$2\frac{2}{3}$ and $2\frac{25}{42} \approx 2.596$, respectively. Let us first
introduce necessary definitions.

\section{Preliminaries}
\begin{definition}[Factor and Periodicity \cite{Breslauer1996}]
A string $x$ is a factor of a string $s$ if $s = x^i y$ for some positive
integer $i$ and a prefix $y$ of $x$ (where $y$ may be empty). The shortest such
factor of a non-empty string $s$, denoted as $\textit{factor}(s)$, defines the
period of $s$, which is denoted as $\textit{period}(s) = |\textit{factor}(s)|$.
\end{definition}

\begin{definition}[Semi-Infinite Strings and Equivalence \cite{Breslauer1996}] \label{def:semi-infinite}
A semi-infinite string $s = a_1 a_2 \dots$ is said to be periodic if $s = xs$
for some non-empty string $x$. The shortest such $x$ is called the factor of
$s$. Two (periodic semi-infinite) strings $s$ and $t$ are equivalent if their
factors are cyclic shifts of each other, i.e., if there are strings $x$ and $y$
such that $\textit{factor}(s) = xy$ and $\textit{factor}(t) = yx$. Otherwise,
they are inequivalent.
\end{definition}

\begin{example}
Let $s = (ab)^\infty = abababab\dots$ and $t = (ba)^\infty = babababa\dots$.
Then $\textit{factor}(s) = ab$ and $\textit{factor}(t) = ba$, and taking $x =
a$, $y = b$ gives $\textit{factor}(s) = xy$ and $\textit{factor}(t) = yx$. Hence
$s$ and $t$ are equivalent.

In contrast, let $s' = (aabb)^\infty = aabbaabb\dots$ and $t' = (abab)^\infty =
abababab\dots$. Since $\textit{factor}(t') = abab$ is not a cyclic shift of
$\textit{factor}(s') = aabb$, the strings $s'$ and $t'$ are inequivalent.
\end{example}

For each string $s$, let $s^i$ denote the concatenation of $i$ consecutive
copies of $s$. Furthermore, let $s^\infty$ denote the semi-infinite string $sss
\dots$, and let $s_\infty = \textit{factor}(s)^\infty$ denote the periodic
semi-infinite string that begins with $s$. Note that, in general, $s^\infty \neq
s_\infty$.

\begin{example}
Let $s = aaba$. Then $\textit{factor}(s) = aab$, and we have
\[
s^\infty = aabaaaba\dots \qquad \text{while} \qquad s_\infty =
\textit{factor}(s)^\infty = aabaabaab\dots,
\]
giving that $s^\infty \neq s_\infty$.
\end{example}

We extend the notion of equivalence from Definition~\ref{def:semi-infinite} to
finite strings: a finite string $s$ is equivalent to a (periodic semi-infinite)
string $t$ if $s_\infty$ and $t$ are equivalent, and two finite strings $s$ and
$s'$ are equivalent if $s_\infty$ and $s'_\infty$ are equivalent, i.e., if
$\textit{factor}(s)$ and $\textit{factor}(s')$ are cyclic shifts of each other.
In particular, every string $s$ is equivalent to $s_\infty$.

\begin{lemma} \label{lemma:string_containment}
Suppose that a string $x$ is contained in a string $y$, and $y$ is contained in
a string $z$. If $x$ and $z$ are equivalent, then $y$ is also equivalent to $x$
and $z$.
\end{lemma}

\subsection{Cycles and Periodicity}
The following facts establish the connection between a cycle in the distance
graph $G_S$ and the periodicity of the strings obtained by breaking that cycle.
Let $C = s_{i_1}, \dots, s_{i_r}, s_{i_1}$ be a cycle in $G_S$. Without loss of
generality, assume that $C$ has the minimum weight among all cycles in $G_S$
containing $s_{i_1}, \dots, s_{i_r}$. Let $w(C)$ denote the weight of this
cycle. We define $\langle s_{i_1}, \dots, s_{i_r} \rangle$ as the string
obtained by breaking the cycle $C$ at $s_{i_r}$.

\begin{Fact} \label{lemma:cycle_substring}
The string $\langle s_{i_1}, \dots, s_{i_r} \rangle$ is a substring of
$\textit{factor}(\langle s_{i_1}, \dots, s_{i_r} \rangle)s_{i_1} = \langle
s_{i_1}, \dots, s_{i_r}, s_{i_1} \rangle$.
\end{Fact}

\begin{Fact} \label{lemma:cycle_factor}
The factor of the broken cycle string can be expressed as:
\begin{equation}
\textit{factor}(\langle s_{i_1}, \dots, s_{i_r} \rangle) = \textit{pref}(s_{i_1}, s_{i_2}) \dots \textit{pref}(s_{i_{r-1}}, s_{i_r}) \textit{pref}(s_{i_r}, s_{i_1})
\end{equation}
\end{Fact}

From Lemma \ref{lemma:cycle_factor}, three important corollaries follow
immediately:

\begin{corollary} \label{cor:cycle_weight_period}
The weight of the cycle equals the period of the generated superstring:
\begin{equation}
w(C) = d(s_{i_1}, s_{i_2}) + \dots + d(s_{i_{r-1}}, s_{i_r}) + d(s_{i_r}, s_{i_1}) = \textit{period}(\langle s_{i_1}, \dots, s_{i_r} \rangle)
\end{equation}
\end{corollary}

\begin{corollary} \label{cor:cycle_equivalence}
The strings generated by different breaking points of the cycle, namely $\langle
s_{i_1}, \dots, s_{i_r} \rangle, \dots, \langle s_{i_r}, s_{i_1}, \dots,
s_{i_{r-1}} \rangle$, are all mutually equivalent.
\end{corollary}

\begin{corollary} \label{cor:cycle_extension_equivalence} It holds that $\textit{factor}(\langle s_{i_1}, \dots, s_{i_r}, s_{i_1} \rangle) = \textit{factor}(\langle s_{i_1}, \dots, s_{i_r} \rangle)$. That is, the string $\langle s_{i_1}, \dots, s_{i_r}, s_{i_1} \rangle$ is equivalent to $\langle s_{i_1}, \dots, s_{i_r} \rangle$.
\end{corollary}

It is important to note that, in general, the individual strings $s_{i_1},
\dots, s_{i_r}$ are neither mutually equivalent, nor is $s_{i_1}$ equivalent to
the combined string $\langle s_{i_1}, \dots, s_{i_r} \rangle$.

Finally, if two cycles generate equivalent strings, they can be merged without
increasing the cycle weight:

\begin{lemma} \label{lemma:cycle_merging}
Let $C = s_{i_1}, \dots, s_{i_r}, s_{i_1}$ and $C' = s_{j_1}, \dots, s_{j_l},
s_{j_1}$ be two distinct cycles. If the string $\langle s_{i_1}, \dots, s_{i_r}
\rangle$ is equivalent to $\langle s_{j_1}, \dots, s_{j_l} \rangle$, then there
exists a third cycle $\tilde{C}$ with weight $w(C)$ containing all vertices from
both $C$ and $C'$.
\end{lemma}
